\documentclass[tikz,border=18pt]{standalone}
\usepackage{luatexja}
\usepackage{tikz}
\usetikzlibrary{shapes.geometric,arrows.meta,positioning,calc}

\definecolor{secblue}{RGB}{20, 70, 130}
\definecolor{lightblue}{RGB}{236, 246, 255}
\definecolor{secgreen}{RGB}{25, 125, 45}
\definecolor{lightgreen}{RGB}{236, 255, 240}
\definecolor{secorange}{RGB}{210, 95, 20}
\definecolor{lightorange}{RGB}{255, 246, 230}
\definecolor{secred}{RGB}{185, 30, 30}
\definecolor{lightred}{RGB}{255, 238, 238}

\begin{document}
\begin{tikzpicture}[>=Stealth, font=\small]

  % ---------------------------------------------------------------
  % 上段：認証の3大要素（知識・所持・生体）
  % 横幅 5.8cm、間隔 6.4cm と十分なゆとりを確保
  % ---------------------------------------------------------------
  \node[draw=secblue, fill=lightblue, rounded corners=8pt, line width=1.5pt, minimum width=5.8cm, minimum height=2.4cm, align=center] (know) at (0, 5.8) {
    \textbf{\large 【知識情報】}\\[4pt]
    本人だけが「知っていること」\\[4pt]
    \footnotesize パスワード、暗証番号 (PIN)、合言葉
  };

  \node[draw=secgreen, fill=lightgreen, rounded corners=8pt, line width=1.5pt, minimum width=5.8cm, minimum height=2.4cm, align=center] (poss) at (6.5, 5.8) {
    \textbf{\large 【所持情報】}\\[4pt]
    本人だけが「持っているもの」\\[4pt]
    \footnotesize 認証アプリ(OTP)、SMS通知、ICカード
  };

  \node[draw=secorange, fill=lightorange, rounded corners=8pt, line width=1.5pt, minimum width=5.8cm, minimum height=2.4cm, align=center] (bio) at (13.0, 5.8) {
    \textbf{\large 【生体情報】}\\[4pt]
    本人自身の「身体的特徴」\\[4pt]
    \footnotesize 指紋認証、顔認証、虹彩・静脈スキャン
  };

  % ---------------------------------------------------------------
  % 中段：単一要素認証の危険性警告
  % ---------------------------------------------------------------
  \node[draw=secred, fill=lightred, rounded corners=6pt, dashed, line width=1.2pt, align=center, minimum width=18.8cm, inner sep=10pt] (single_risk) at (6.5, 3.0) {
    \textbf{\textcolor{secred}{\large ▲ 単一要素認証（例：パスワードのみ）の脆弱性とリスク}}\\[4pt]
    \footnotesize 「知識」が1つ漏洩（使い回しやリスト型攻撃・総当たり攻撃）しただけで、第三者にアカウントが即座に乗っ取られる！
  };

  % ---------------------------------------------------------------
  % 下段：多要素認証（MFA）のアーキテクチャ
  % ---------------------------------------------------------------
  \node[draw=secblue, fill=secblue!12, rounded corners=10pt, line width=1.8pt, minimum width=18.8cm, minimum height=2.8cm, align=center] (mfa) at (6.5, -0.6) {
    \textbf{\Large 多要素認証（MFA: Multi-Factor Authentication）}\\[6pt]
    \textbf{\normalsize 異なる２つ以上の「独立した要素」を組み合わせて本人を確認}\\[4pt]
    \small 【最強の組み合わせ例】「知識（パスワード）」＋「所持（手元のスマホ認証アプリ）」\\[3pt]
    \footnotesize \textcolor{secblue}{\textbf{⇒ たとえパスワードが漏洩しても、手元の物理スマートフォンがないため不正ログインを確実に阻止！}}
  };

  % ---------------------------------------------------------------
  % 矢印接続：上段3要素から下段MFAへ接続
  % ---------------------------------------------------------------
  \draw[->, very thick, secblue] (know.south) -- (know.south |- single_risk.north);
  \draw[->, very thick, secblue] (know.south |- single_risk.south) -- (know.south |- mfa.north);

  \draw[->, very thick, secgreen] (poss.south) -- (single_risk.north);
  \draw[->, very thick, secgreen] (single_risk.south) -- (mfa.north);

  \draw[->, very thick, secorange] (bio.south) -- (bio.south |- single_risk.north);
  \draw[->, very thick, secorange] (bio.south |- single_risk.south) -- (bio.south |- mfa.north);

\end{tikzpicture}
\end{document}
